\makeabstract
{
Electrodeposited Bimetallic Catalysts for Impurity-Tolerant Electrocatalytic Hydrogenation of Muconic Acid to Adipic Acid
}
{
Richard Lin (1), Abriana Ferrari (2), Devanshi Mistry (3,4), and Jean-Philippe Tessonnier (3,4)
}
{
\textsuperscript{1}Department of Chemistry, Amherst College, Amherst, MA, USA\\
\textsuperscript{2}Department of Biology, Lebanon Valley College, Annville, PA, USA\\
\textsuperscript{3}Center for Biorenewable Chemicals (CBiRC), Ames, IA, USA\\
\textsuperscript{4}Dept. of Chemical and Biological Engineering, Iowa State University, Ames, IA, USA
}
{
The application of electrocatalysis to biomass conversion, biobased materials, and bioenergy is a promising approach for transitioning away from a fossil fuel-dependent economy. Of particular interest is the synthesis of adipic acid, an industrially valuable compound and a key monomer of nylon 6,6. This study examines the electrocatalytic hydrogenation (ECH) of fermentation-derived cis,cis-muconic acid (ccMA) to adipic acid (AA), with an emphasis on tuning the catalyst composition to enable efficient conversion directly in unpurified fermentation broth. This integrated approach has the potential to improve overall process efficiency by eliminating the need for downstream purification steps to extract ccMA before electrocatalytic conversion. A key challenge in this approach is catalyst deactivation by biomolecules, particularly sulfur-containing amino acids, that are inherently present in the fermentation broth. Previous work in our group has shown that poisoned catalysts can be regenerated in the absence of these impurities using electrochemical techniques like cyclic voltammetry (CV) and reductive stripping. Therefore, developing robust and impurity-tolerant catalysts is crucial for enabling the direct use of fermentation broth as the reaction medium.
}
{
One potential strategy to improve sulfur tolerance is by modifying the electronic structure of the catalyst through alloying. Electrodeposition provides an effective and low-cost method for synthesizing such alloyed catalysts, offering uniform deposition and simple implementation. In this study, PdAu, PdRu, PdBi, and PdPt bimetallic catalysts were deposited on porous nickel foam using chronoamperometry (CA). The resulting catalysts were tested for the electrocatalytic hydrogenation of ccMA to AA, and quantitative 1H-NMR spectroscopy was used to characterize the production of adipic acid.
}
{
The PdAu/Ni catalyst outperformed all other catalysts, reaching near full conversion of ccMA after 3 hours. Scanning electron microscopy (SEM) images of the catalysts showed a wide distribution in size and morphology of electrodeposited metals, while energy dispersive X-ray spectroscopy (EDS) and X-ray diffraction (XRD) were performed on the PdAu/Ni catalyst to confirm deposition of palladium and gold over the nickel foam surface and alloying of palladium and gold.
}
{
A set of metal and metal alloy electrocatalysts synthesized through electrodeposition were screened to investigate their effectivity in converting ccMA to AA. Future directions include testing catalysts for poisoning resistance. Investigation of the electronic properties of alloys through density functional theory (DFT) and other computational methods may suggest other bimetallic catalysts to be experimentally tested in the lab.
}

\newpage
